







\section{Method}
\label{sec:method}



\subsection{Problem setup}
\label{sec:method_problem}
Given a set of $V$ input images $\{\im_i\}_{i=1}^{V}$ with $\im_i \in \real^{H \times W \times 3}$, our goal is to recover the underlying 3D scene geometry, expressed in the coordinate frame of the first view. We adopt a depth-ray representation: for each image $\im_i$, the model predicts a per-pixel depth map $\depthmap_i \in \real^{H \times W}$ and a dense ray map $\raymap_i \in \real^{H \times W \times 6}$.
Each pixel of the ray map encodes a 3D origin $\raymap^o \in \real^3$ and an unnormalized direction $\raymap^d \in \real^3$, so that a 3D point in world coordinates is obtained as $\mathbf{X} = \raymap^o + \depthmap(u,v) \cdot \raymap^d$.
Per-view camera-to-world rotations $\mathbf{R}_i \in SO(3)$, translations $\mathbf{t}_i \in \real^3$, and intrinsic matrices $\mathbf{K}_i \in \real^{3 \times 3}$ are recovered from the predicted ray maps following~\citet{lin2025depth}.


\subsection{Hypothesis}
\label{sec:hypothesis}

Two recent analyses motivate our approach.
First, \citet{jacobs2025raptor} show via post-hoc distillation that the $L$ layers of a trained ViT can be accurately approximated by a few looped blocks.
Second, \citet{stary2025understanding} probe DUSt3R~\citep{dust3r} layer by layer and show that its predicted pointmaps progressively refine across decoder depth, revealing iterative refinement of geometry inside multi-view transformers even though their layers do not share weights.

We hypothesize that this structure can be made explicit by applying a looped block to an evolving state, and that the resulting recurrence performs \emph{directional refinement} of the state, where the direction of $\z_k$ converges to the endpoint direction over the trained step range
(\cref{sec:method_recurrence,sec:analysis}).
In contrast to task-space iterative refinement methods such as RAFT~\citep{RAFT,lipson2021raft}, which decode after every step and apply a sequence loss on the output, we refine an internal state with a looped block and supervise only at the final step. This avoids running the decoder and computing a loss at every intermediate step, sparing $(K{-}1)$ decoder forward and backward passes per training iteration.

We model the recurrence as a time-conditioned discrete update over the partition $0 = t_0 < t_1 < \cdots < t_K = 1$ of the unit interval:
\begin{equation}
    \z_{k+1} = f_\theta(\z_k,\, t_k,\, t_{k+1})\,,
\label{eq:discrete_recurrence}
\end{equation}
where $f_\theta$ is a looped block conditioned on the continuous time interval $(t_k, t_{k+1})$.
Conditioning on continuous time, rather than on a discrete step index as in prior weight-tied transformers~\citep{dehghani2019universal}, decouples the block from any specific value of $K$ and lets a single set of weights cover a range of step counts at inference.
\Cref{sec:method_recurrence} instantiates $f_\theta$ as a shared transformer block with frame and global attention.

\begin{figure*}[t]
  \centering
  \includegraphics[width=\linewidth]{figures/method_overview_v3.pdf}
    \caption{\textbf{Method overview.} $V$ input images are encoded by a shared DINOv2~\citep{oquab2023dinov2} backbone. A single looped transformer block with frame-wise and global attention sub-blocks is then applied recurrently to the resulting tokens for $K$ steps, with $K$ sampled per batch from $[K_\text{min}, K_\text{max}]$ during training. Two heads decode the final tokens into per-view depth and ray predictions.}
  \label{fig:method_overview}
  \vspace{-4mm}
\end{figure*}




\subsection{Architecture}
\label{sec:method_arch}

\parsection{Patch encoder and tokens.}
We initialize the per-view state $\z_0$ from a pretrained DINOv2~\citep{oquab2023dinov2} encoder, which maps each input image to an $\frac{H}{P} \times \frac{W}{P}$ grid of patch tokens.
We prepend a per-view copy of $R$ learnable register tokens~\citep{darcet2024vision} and of a learnable camera token to each view's token sequence, with the underlying parameters tied across views. The camera token uses two parameter sets: one for the reference (first) view, and one tied across all other views.
We encode patch positions with 2D rotary position embeddings~\citep{heo2024ropevit} and assign special tokens a sentinel position outside the patch grid.
We then apply a looped transformer block $K$ times to $\z_0$, with $K$ randomly sampled per training batch (\cref{sec:method_recurrence}), yielding the final state $\z_K$ passed to the decoder heads.

\parsection{Decoder heads.}
We pass $\z_K$ through two parallel decoder branches, each comprising a shallow transformer followed by an output head.
Each decoder transformer uses the same pre-norm $\text{Attn}+\text{MLP}$ block design as the shared recurrent block (\cref{sec:method_recurrence}), but without LayerScale.
The ray decoder uses a linear pixel-shuffle head to produce the per-pixel ray map $\raymap_\theta \in \real^{H \times W \times 6}$.
The depth decoder uses a convolutional depth head following~\citet{wang2025moge} to avoid block artifacts at patch boundaries (\cref{sec:suppl_two_stage}), and produces the depth map $\depthmap_\theta \in \real^{H \times W}$ and a depth confidence map $c_\depthmap \in \real^{H \times W}$.
Following \citet{lin2025depth}, we additionally include a camera MLP head that decodes a tuple $\mathbf{c}_\theta = (\mathbf{t}_\theta, \mathbf{q}_\theta, \mathbf{f}_\theta) \in \real^3 \times \mathbb{S}^3 \times \real^2$, comprising translation, unit rotation quaternion, and field of view, from the per-view camera tokens at the output of the ray decoder (\cref{sec:method_training}). It provides a faster alternative to the rays-derived recovery of \cref{sec:method_problem}, which remains our default at inference.
We obtain world points analytically as $\mathbf{X}_\theta = \raymap_\theta^o + \depthmap_\theta \cdot \raymap_\theta^d$.

\begin{figure}[t]
    \centering
    \begin{subfigure}[t]{\linewidth}
        \centering
        \includegraphics[width=0.9\linewidth]{figures/overall_convergence_metrics.png}
        \caption{%
            \textbf{Task quality across recurrent iterations.}
            Decoding the residual stream $\z_k$ at every iteration $k\!\in\!\{1,\dots,16\}$ yields monotone improvement of pose and pointmap metrics across the trained step range.%
        }
        \label{fig:refinement-metrics}
    \end{subfigure}
    \begin{subfigure}[t]{\linewidth}
        \centering
        \includegraphics[width=0.9\linewidth]{figures/overall_convergence.pdf}
        \caption{%
            \textbf{Residual-stream convergence.}
            \emph{Left:} cosine similarity between $\z_k$ and the final $\z_{16}$. \emph{Middle:} relative update norm $\lVert \Delta \z_k \rVert / \lVert \z_k \rVert$. \emph{Right:} feature norm $\lVert \z_k \rVert_{2}$.%
        }
        \label{fig:refinement-convergence}
    \end{subfigure}
    \caption{%
        \textbf{Iterative refinement of the residual stream.}
        Per-iteration analysis of \methodname's recurrent block, averaged across the five benchmarks of \cref{tab:comparison_sota_pmap,tab:comparison_sota_pose}.
        Task quality improves monotonically with the iteration count \textbf{(top)}. The recurrence does not contract to a fixed point in feature space (the state norm grows) but its direction stabilizes, with cosine similarity to the final state approaching $1$ and the relative update norm decaying by roughly $5{\times}$ \textbf{(bottom)}. The decoder's input LayerNorm absorbs the norm growth, so the decoded representation effectively converges in direction.%
    }
    \label{fig:refinement}
    \vspace{-4mm}
\end{figure}



\subsection{Looped Block}
\label{sec:method_recurrence}

\parsection{Block design.}
The looped block consists of two attention sub-blocks applied in sequence, following the alternating frame/global design of VGGT~\citep{vggt}.
The first is a frame attention that processes each view independently with 2D rotary position embeddings.
The second is a global attention that operates over the joint sequence of all tokens across all views.
Each sub-block uses a standard pre-norm $\text{Attn}+\text{MLP}$ design with LayerScale~\citep{touvron2021going}.

We condition the block on the time interval $(t_k, t_{k+1})$.
Three channel-wise scale vectors $(\mathbf{s}_\text{attn}, \mathbf{s}_\text{mlp}, \mathbf{s}_\text{out})$ control the block update:
\begin{equation}
\begin{aligned}
    & \z' && = \z_k + \mathbf{s}_\text{attn} \odot \text{LS}_1(\text{Attn}(\text{LN}_1(\z_k)))\,, \\
    & \z'' && = \z' + \mathbf{s}_\text{mlp} \odot \text{LS}_2(\text{MLP}(\text{LN}_2(\z')))\,, \\
    & \z_{k+1} && = \mathbf{s}_\text{out} \odot \z''\,,
\end{aligned}
\label{eq:gated_block}
\end{equation}
where $\text{LN}$ is layer normalization, $\text{LS}$ is LayerScale, and $\odot$ is channel-wise multiplication broadcast over the sequence dimension.
We compute the scales via a zero-initialized MLP such that $\mathbf{s} = \mathbf{1} + \text{MLP}(\gamma(t_k, t_{k+1}))$, where $\gamma$ concatenates the sinusoidal embeddings of $t_k$ and $t_{k+1}$.
We ablate simpler variants of this block (no $\mathbf{s}_\text{out}$, no time conditioning, and untied per-step blocks) in \cref{sec:analysis}.



\parsection{Variable step count.}
We train a single set of weights to serve as a 
$K$-elastic
family: the same block supports a range of step counts $K$ at inference, exposed as a compute knob.
Specifically, we sample $K \sim \text{Beta}(\alpha, \beta)$ per batch, scaled and rounded into $[K_\text{min}, K_\text{max}]$, and apply the block $K$ times along the uniform partition $0 = t_0 < t_1 < \cdots < t_K = 1$ with $t_k = k/K$,
where the application from $\z_k$ to $\z_{k+1}$ is conditioned on the interval $(t_k, t_{k+1})$.
At inference, we run the block on a uniform grid of $K_\text{inf}$ steps; varying $K_\text{inf}$ trades compute for accuracy within the trained range $[K_\text{min}, K_\text{max}]$, with degradation observed when $K_\text{inf}$ is pushed substantially outside it (\cref{sec:analysis}, \cref{sec:beyond_kmax}).


\parsection{Directional refinement.}
Unlike deep equilibrium networks~\citep{bai2019deep}, our recurrence does not converge to a fixed point in feature space. Instead, the state norm $\|\z_k\|$ grows monotonically with $k$.
However, two empirical signatures characterize its dynamics within the trained step range (\cref{fig:refinement}).
First, $\cos(\z_k, \z_K)$ rises monotonically toward $1$ as $k \to K$, which means that each step moves the state closer in direction to the endpoint.
Second, the relative update magnitude $\|\Delta \z_k\|/\|\z_k\|$ decays from $\sim 0.5$ at the first step to $\sim 0.1$ at the last, indicating a genuine slowdown of motion rather than a constant rescaling.
Because each decoder branch is pre-norm (\cref{sec:method_arch}), the LayerNorm at the start of its first transformer block absorbs the component of $\Delta \z_k$ parallel to $\z_k$ that drives the norm growth, and the decoded representation effectively converges in direction. 
We refer to this behavior as \emph{directional refinement}, distinct from RAFT-style task-space refinement~\citep{RAFT} that operates on the output and contracts in absolute magnitude.




\subsection{Training}
\label{sec:method_training}

\parsection{Losses.}
We supervise the model directly on the predicted geometry with five complementary
loss terms
covering depth, rays, world points, and camera parameters.
Following DUSt3R~\citep{dust3r}, predictions and ground truth are independently normalized prior to loss computation: given valid 3D points $\{\mathbf{X}_j\}_{j \in \Omega}$, we define the inverse normalization scale
\begin{equation}
    s = \left(\frac{1}{|\Omega|} \sum_{j \in \Omega} \lVert \mathbf{X}_j \rVert_2\right)^{\!-1}\,,
\label{eq:normalization_scale}
\end{equation}
computed separately for the predicted ($\hat{s}$) and ground-truth ($\bar{s}$) point clouds.
The per-sample training loss is then:
\begin{equation}
\begin{split}
    \loss ={} & \lVert \hat{s}\,\depthmap_\theta - \bar{s}\,\depthmap \rVert_2 + \loss_\text{grad}(\hat{s}\,\depthmap_\theta, \bar{s}\,\depthmap) + \lVert \hat{s}\,\raymap_\theta - \bar{s}\,\raymap \rVert_1 \\
    & + \lVert \hat{s}\,\mathbf{X}_\theta - \bar{s}\,\mathbf{X} \rVert_2 + \loss_\text{cam}(\mathbf{c}_\theta, \mathbf{c})\,,
\end{split}
\label{eq:total_loss}
\end{equation}
where $\mathbf{X}_\theta = \raymap_\theta^o + \depthmap_\theta \cdot \raymap_\theta^d$ is the analytically derived predicted point cloud, $\loss_\text{grad}$ is a multi-scale $\ell_1$ loss on horizontal and vertical depth gradients~\citep{Hu2019RevisitingSI, lin2025depth}, and $\loss_\text{cam}$ decomposes into separately weighted $\ell_1$ terms on the translation, rotation, and field-of-view components of $\mathbf{c}_\theta$.
The pointmap term ties the depth and ray heads through a joint geometric signal. %

\begin{table}[t]
    \caption{%
        \textbf{Pointmap accuracy.} Relative $\ell_2$ distance (\textbf{Rel.~L2}~$\downarrow$) and inlier ratio (\textbf{IR}~$\uparrow$) on the global pointmap after a Sim(3) alignment to the ground truth, across five benchmarks. The \colorbox{tabred}{best}, \colorbox{taborange}{second}, and \colorbox{tabyellow}{third} ranked results are highlighted. VGGT-$\Omega$~\citep{vggt_omega} is concurrent work, reported here for completeness.
    }
    \label{tab:comparison_sota_pmap}
    \footnotesize
    \setlength{\tabcolsep}{4pt}
    \centering
    \noindent\makebox[0.85\linewidth][c]{%
    \begin{tabular}{l rr rr rr rr rr}
        \toprule
        \multirow{2}{*}[-4pt]{\textbf{Method}}
            & \multicolumn{2}{c}{\textbf{DTU}} & \multicolumn{2}{c}{\textbf{ETH3D}} & \multicolumn{2}{c}{\textbf{7-Scenes}} & \multicolumn{2}{c}{\textbf{ScanNet++}} & \multicolumn{2}{c}{\textbf{nuScenes}} \\
        \cmidrule(lr){2-3}\cmidrule(lr){4-5}\cmidrule(lr){6-7}\cmidrule(lr){8-9}\cmidrule(lr){10-11}
            & Rel.~L2 & IR & Rel.~L2 & IR & Rel.~L2 & IR & Rel.~L2 & IR & Rel.~L2 & IR \\
        \midrule
        MASt3R~\citep{leroy2025grounding}                       & \cellcolor{tabyellow}0.011 & 94.9 & 0.340 & 29.1 & 0.076 & 41.7 & 0.251 & 14.5 & 0.360 & 11.4 \\
        MASt3R-SfM~\citep{duisterhof2024mast3rsfm}              & \cellcolor{tabred}0.009 & 96.9 & 0.095 & 54.1 & 0.051 & 64.7 & 0.042 & 69.7 & 0.311 & 18.4 \\
        MapAnything~\citep{keetha2026mapanything}               & 0.014 & 95.2 & 0.227 & 40.6 & 0.044 & 67.9 & \cellcolor{tabyellow}0.019 & \cellcolor{tabyellow}89.2 & 0.089 & \cellcolor{tabyellow}51.9 \\
        Pi3~\citep{wang2025pi}                                  & \cellcolor{tabred}0.009 & \cellcolor{tabred}97.3 & \cellcolor{tabyellow}0.034 & \cellcolor{tabyellow}66.8 & \cellcolor{tabred}0.032 & \cellcolor{tabred}77.8 & \cellcolor{tabred}0.014 & \cellcolor{tabred}94.3 & \cellcolor{tabyellow}0.078 & 51.0 \\
        VGGT~\citep{vggt}                                       & \cellcolor{taborange}0.010 & 95.8 & 0.053 & 52.6 & 0.042 & 70.8 & 0.034 & 68.4 & 0.081 & 42.3 \\
        VGGT-$\Omega$-1B~\citep{vggt_omega}                        & \cellcolor{tabred}0.009 & \cellcolor{taborange}97.2 & \cellcolor{tabred}0.024 & \cellcolor{tabred}78.6 & 0.039 & 65.3 & 0.032 & 70.6 & \cellcolor{tabred}0.055 & \cellcolor{tabred}62.3 \\
        DA3-L~\citep{lin2025depth}                              & \cellcolor{taborange}0.010 & \cellcolor{tabyellow}97.1 & 0.211 & 49.9 & 0.039 & 69.6 & 0.051 & 48.1 & 0.141 & 27.0 \\
        DA3-G~\citep{lin2025depth}                              & \cellcolor{taborange}0.010 & \cellcolor{tabyellow}97.1 & 0.129 & 64.7 & \cellcolor{tabyellow}0.037 & \cellcolor{tabyellow}71.1 & 0.041 & 57.8 & 0.080 & 42.0 \\
        \midrule
        \textbf{Ours}                                           & \cellcolor{tabred}0.009 & \cellcolor{tabyellow}97.1 & \cellcolor{taborange}0.026 & \cellcolor{taborange}78.3 & \cellcolor{taborange}0.035 & \cellcolor{taborange}74.2 & \cellcolor{taborange}0.015 & \cellcolor{taborange}93.3 & \cellcolor{taborange}0.067 & \cellcolor{taborange}58.5 \\
        \bottomrule
    \end{tabular}}
\end{table}



\begin{table}[t]
    \caption{%
        \textbf{Camera pose accuracy.} Area under the cumulative pose-error curve at $3^{\circ}$ (\textbf{AUC@3}~$\uparrow$) and $30^{\circ}$ (\textbf{AUC@30}~$\uparrow$), across five benchmarks.         VGGT-$\Omega$~\citep{vggt_omega} is concurrent work, reported here for completeness. \methodname{} ranks first or second on nine of ten cells and is in the top three on every cell.
    }
    \label{tab:comparison_sota_pose}
    \footnotesize
    \setlength{\tabcolsep}{6pt}
    \centering
    \noindent\makebox[0.85\linewidth][c]{%
    \begin{tabular}{l rr rr rr rr rr}
        \toprule
        \multirow{2}{*}[-4pt]{\textbf{Method}}
            & \multicolumn{2}{c}{\textbf{DTU}} & \multicolumn{2}{c}{\textbf{ETH3D}} & \multicolumn{2}{c}{\textbf{7-Scenes}} & \multicolumn{2}{c}{\textbf{ScanNet++}} & \multicolumn{2}{c}{\textbf{nuScenes}} \\
        \cmidrule(lr){2-3}\cmidrule(lr){4-5}\cmidrule(lr){6-7}\cmidrule(lr){8-9}\cmidrule(lr){10-11}
        \multicolumn{1}{r}{\normalfont AUC}
            & @3 & @30 & @3 & @30 & @3 & @30 & @3 & @30 & @3 & @30 \\
        \midrule
        MASt3R~\citep{leroy2025grounding}                       & 21.6 & 81.7 & 35.2 & 57.6 & 9.3 & 70.6 & 15.5 & 43.9 & 9.4 & 62.7 \\
        MASt3R-SfM~\citep{duisterhof2024mast3rsfm}              & 40.2 & 91.4 & 52.8 & \cellcolor{taborange}85.0 & \cellcolor{taborange}16.4 & 80.3 & 38.7 & 85.0 & 9.1 & 70.9 \\
        MapAnything~\citep{keetha2026mapanything}               & 18.1 & 88.8 & 37.1 & 73.8 & 8.2 & 74.6 & \cellcolor{tabyellow}70.6 & \cellcolor{tabyellow}96.7 & 37.8 & 82.8 \\
        Pi3~\citep{wang2025pi}                                  & 70.2 & 97.3 & 43.2 & \cellcolor{tabyellow}84.9 & 11.3 & \cellcolor{tabyellow}81.3 & \cellcolor{taborange}76.5 & \cellcolor{taborange}97.3 & 24.8 & 82.4 \\
        VGGT~\citep{vggt}                                       & \cellcolor{tabred}96.5 & \cellcolor{tabred}99.8 & 35.4 & 82.8 & 11.1 & 79.1 & 15.1 & 74.7 & \cellcolor{tabyellow}39.2 & \cellcolor{tabyellow}83.8 \\
        VGGT-$\Omega$-1B~\citep{vggt_omega}                        & \cellcolor{tabyellow}77.4 & \cellcolor{tabyellow}98.1 & \cellcolor{taborange}64.1 & \cellcolor{tabred}95.4 & \cellcolor{tabred}21.3 & \cellcolor{tabred}87.0 & 29.9 & 87.3 & \cellcolor{taborange}42.4 & \cellcolor{tabred}85.5 \\
        DA3-L~\citep{lin2025depth}                              & 69.2 & 97.3 & 38.8 & 75.3 & 11.8 & 79.7 & 7.6 & 71.8 & 11.0 & 74.7 \\
        DA3-G~\citep{lin2025depth}                              & 74.9 & 97.9 & \cellcolor{tabyellow}55.4 & 82.4 & 12.0 & 80.9 & 26.4 & 80.2 & 37.7 & 82.0 \\
        \midrule
        \textbf{Ours}                                           & \cellcolor{taborange}83.2 & \cellcolor{taborange}98.8 & \cellcolor{tabred}66.0 & \cellcolor{tabred}95.4 & \cellcolor{tabyellow}13.9 & \cellcolor{taborange}81.7 & \cellcolor{tabred}79.4 & \cellcolor{tabred}98.0 & \cellcolor{tabred}43.4 & \cellcolor{taborange}85.3 \\
        \bottomrule
    \end{tabular}}
\end{table}










\parsection{Optimization.}
Training proceeds in two stages.
The first stage trains the model end-to-end with all loss terms, using plain $\ell_2$ regression on the depth and pointmap terms and a linear pixel-shuffle depth head.
The second stage swaps in the convolutional depth head described in \cref{sec:method_arch} and finetunes the depth decoder while freezing all other parameters. The ray and camera losses are disabled in the second stage. The depth term applies DUSt3R-style confidence weighting~\citep{dust3r}, $\loss_\text{unc}(c, \mathbf{a}, \mathbf{b}) = c \lVert \mathbf{a} - \mathbf{b} \rVert_2 - \lambda_c \log c$, parameterized by the predicted per-pixel uncertainty $c_\depthmap$, while the pointmap term remains plain $\ell_2$ on the analytically derived points.













